\chapter{Problem formulation}

A previous thesis \cite{insaghi2025thesis} already developed and analyzed 
a MARL algorithm suitable for this task. It used an environment made of
gaussians, which defined the reward function. However, the real world is 
more complex and made of currents. This thesis, thus, brings these contributions:
\begin{itemize}
    \item A new, more realistic environment, composed of currents, salinity and turbidity fields, created with Oceananigans;
    \item Another dimension to the problem, so that it goes from 2D to 3D;
    \item The use of a simulator, SwarmSwIM, to connect the environment with the agents and the MARL algorithm.
\end{itemize}

In the following sections, all those components are described in detail, 
including the parameters used to create the environment with the solver
and how agents are commanded via the simulator.

\section{SwarmSwIM Simulator}

SwarmSwIM \cite{swim_simulator} is a simulator developed to help researchers
and pratictioners to evaluate and enhance the control systems and decision-making
capabilities of underwater agents (AUVs). Completely Python-based, it provides
tools for current modeling, agent control and visual representation, balancing 
computational efficiency with sufficient environmental fidelity to facilitate large-scale
swarm simulations. The low-level control of the agents is approximated by the 
closed-loop control response of each agent, such as depth and heading keeping, in the 
resolution of the agent motion. 

Each member of the swarm is an \textit{agent} and is represented by a point
in the 3D space defined by the domain of the environment. The main cycle of it 
is made by simulation and agent ticks, where the former refers to the global update
while the latter to the local update of each agent. 

\section{Agent}

An agent is a point in the 3D space and its state is determined by four DoF: the 3D
position and the yaw angle. Roll and pitch are not considered, and yaw is fixed during 
an episode for simplicity. Agent control is achieved by three schemes:
\begin{enumerate}
    \item \textit{Depth}: Set to heave, such that the vertical velocity is taken as input;
    \begin{equation}
        z_{t+1} = (v_z + \epsilon) \Delta t + z_t \quad \text{with } \epsilon \sim \mathcal{N}(0, 1)
    \end{equation}
    \item \textit{Heading}: Set to yawrate but the yaw is always fixed;
    \begin{equation}
        \psi_{t+1} = (v_\psi + \epsilon) \Delta t + \psi_t 
    \end{equation}
    \item \textit{Planar}: Set to local velocity, such that current position is updated based on a velocity input command expressed in the body-frame reference system.
    \begin{equation}
        \begin{bmatrix}
            x_{t+1} \\
            y_{t+1} 
        \end{bmatrix} = 
        \begin{bmatrix}
            v_x + \epsilon \\
            v_y + \epsilon
        \end{bmatrix} \Delta t + 
        \begin{bmatrix}
            x_t \\
            y_t
        \end{bmatrix} 
    \end{equation}
\end{enumerate}

\section{Environment}

The environment is a 3D domain where the agents operate. To make it more realistic,
I used Oceananigans.jl \cite{Oceananigans-overview-paper-2025}, which is a 
fast, friendly and flexible software package for finite volume simulations of the 
nonhydrostatic and hydrostatic Boussinesq equations on CPUs and GPUs. An environment 
used in an episode is simply a snapshot of one of the different simulations
I ran with Oceananigans. The environment is composed of three fields:
\begin{itemize}
    \item \textit{Currents}
    \item \textit{Salinity}
    \item \textit{Turbidity}
\end{itemize}
The first one is computed by the solver. The second one is computed using
$N$ salinity sources, initialized randomly at "land" positions (i.e., south 
and west walls in this case) and with a random strength. The field is constructed 
by advecting the particles with the current field, corrected by a diffusion 
term. The third is simply an analytical function of depth:
\begin{equation}
    T(z) = \frac{1}{1 + e^{-(z - z_0)}} \quad \text{with } z_0 = 0.5
\end{equation}

The domain is a cube of $1000 \times 1000 \times 100$ meters, discretized
on a $128 \times 128 \times 64$ grid of cells, with topology (Periodic, Periodic, Bounded). The method chosen for the 
simulation is the nonhydrostatic one, which assumes that the vertical
acceleration is not negligible. This method is more accurate than the 
hydrostatic one, but necessitates more computational resources. 30 simulations
were carried out for training and 10 for testing, each one with a different random seed to ensure variability
in the environment. Each simulation created a single NetCDF file, with a snapshot
containing the fields every 12 minutes. The simulation was run for 18 hours,
with the first 12 used as warmup time and so not recorded. 

% Table with the parameters used for the simulation
\begin{table}[t]
\caption{Parameters used in the Oceananigans nonhydrostatic model\label{tab:1}}\vspace{0.25em} % adds some additional space between caption and \toprule, if desired.
\centering{%
\tagpdfsetup{table/header-rows={1}}
\begin{tabular*}{\textwidth}{@{\hspace*{1.5em}}@{\extracolsep{\fill}}ccc!{\hspace*{1.5em}}ccc@{\hspace*{5.0em}}}
\toprule
Parameter &
Value &
Description \\ \midrule
LX & 1000.0 & Domain (x) \\
LY & 1000.0 & Domain (y) \\
LZ & 100.0 & Domain (z) \\
NX & 128 & Grid points (x) \\
NY & 128 & Grid points (y) \\
NZ & 64 & Grid points (z) \\
$T_{surface}$ & 22.0 / 33.0 & Surface temperature (winter/summer) [°C]\\
$T_{bottom}$ & 20.0 / 22.0 & Bottom temperature (winter/summer) [°C]\\
$S_{baseline}$ & 40.0 / 42.0 & Baseline salinity (winter/summer) [PSU]\\
$V_{wind}$ & 5.0 / 8.0 & Wind speed (winter/summer) \\
$\rho_{air}$ & 1.225 & Air density [$kg/m^3$]\\
$c_{drag}$ & $1.3e^{-3}$ & Drag coefficient \\
$\rho_{water}$ & 1027.0 & Water density [$kg/m^3$]\\
Latitude & 24.5 & Abu Dhabi latitude (degrees N)\\
$\omega_{earth}$ & $7.292e^{-5}$ & Earth rotation rate [$rad/s$]\\
$f_{coriolis}$ & $2\omega_{earth}\sin(\text{Latitude})$ & Coriolis parameter [$rad/s$]\\
$\alpha_t$ & $1.67e^{-4}$ & Thermal expansion coefficient [$K^{-1}$]\\
$\beta_s$ & $0$ & Salinity contraction coefficient [$psu^{-1}$]\\
$\sigma_h$ & 300.0 & Salinity horizontal std [m]\\
$\sigma_v$ & 40.0 & Salinity vertical std [m]\\
Salinity excess & 10.0 & Excess over baseline [PSU]\\
Salinity decay time & 1.0 & Time to relax the field to a gaussian\\
Warmup time & 12.0 & Time to reach a steady state [h]\\
Simulation time & 6.0 & Total simulation time [h]\\
Output frequency & 12.0 & Frequency of output [min]\\
\bottomrule
\end{tabular*}
}%
\end{table}

\subsection{Currents}

Currents are composed by three components: the horizontal velocities ($u,v$),
and the vertical one ($w$). They are prognostic variables integrated 
forward in time by the nonhydrostatic Oceananigans model, which solves the rotating, 
Boussinesq, nonhydrostatic Navier–Stokes equations with an LES subgrid closure. 
The developed script only sets up the forcing, the boundary conditions, the rotation, 
the stratification, and the time-stepping; the flow itself emerges from the solve.
The model integrates the following equation:

\begin{equation}
\frac{\partial \mathbf{\vec{u}}}{\partial t} + \underbrace{(\mathbf{\vec{u}}\cdot\nabla)\mathbf{\vec{u}}}_{\text{advection}} + \underbrace{f \hat{z}\times\mathbf{\vec{u}}}_{\text{coriolis}} = -\underbrace{\nabla p}_{\text{pressure}} + \underbrace{b \hat{z}}_{\text{buoyancy}} + \underbrace{\nabla\cdot\boldsymbol{\tau}_{\rm sgs}}_{\text{subgrid stress}} + \underbrace{\mathbf{F}}_{\text{body force}}
\end{equation}

\begin{equation*}
    \text{with } \nabla\cdot\mathbf{\vec{u}} = 0
\end{equation*}

The fourth term $(-\nabla p)$ is the pressure gradient force. It is a global 
operator (i.e., modifies the entire domain) and essentially is a vector
pointing in the direction of steepest increase in pressure. Because the model 
is Boussinesq, $p$ here is the kinemaic pressure, already divided by $\rho_0 = 1027.0$. 
Moreover, it has no equation of state, it is a Lagrange multiplier that enforces incompressibility, 
and since there is the constraint $\nabla\cdot\mathbf{\vec{u}} = 0$, it is whatever force field 
required to make that true. Mechanically, Oceananigans does thig with a projection
method each timestep:
\begin{itemize}
    \item advances velocity ignoring pressure
    \item solves a Poisson equation for pressure
    \begin{equation}
        \nabla^2 p = \nabla\cdot \mathbf{\vec{u}}^* / \Delta t
    \end{equation}
    \item corrects the velocity with the pressure gradient
    \begin{equation}
        \vec{u} = \vec{u}^* - \Delta t \cdot \nabla p
    \end{equation}
\end{itemize}
The second step is an elliptic solve that uses an FFT along the periodi axis and transforms along the bounded axes.

The fifth term ($b\hat{z}$) is the buoyancy force, TODO.

The sixth term ($\nabla\cdot\boldsymbol{\tau}_{\rm sgs}$) is the subgrid stress, which is a closure for the unresolved scales of motion. 
Essentially, LES applies a spatial filter (the grid) to the Navier-Stokes equations, splitting the velocity into a resolved and an unresolved part.
Physically, $\tau_{\rm sgs}$ is the momentum flux carried by eddies smaller than a grid cell, since in the current model the grid cells cannot see them. 

The last term ($\mathbf{F}$) embodies every forcing that is not in the standard list of fixed physics solved by Oceananigans.
It is a body force, applied throughout the volume at every grid cell. 

Wind speed enters whrough the bulk drag law:
\begin{equation}
\tau = \rho_{\rm air},C_d,U_{10}^2 \quad\Longrightarrow\quad \tau_{\rm kin} = \frac{\rho_{\rm air} C_d U_{10}^2}{\rho_{\rm water}}
\end{equation}

And that is applied as a flux boundary condition at the surface. Momentum enters through the surface and
is transferred downward by turbulence. 
Surface stress plus rotation gives the classical Ekman spiral: the current at the surface
is deflected to the right of the wind (northern hemisphere) and rotates further right with depth while decaying. 

Only temperature is buoyancy-active: salinity exerts zero feedback on momentum, it is advected and LES-mixed by the currents but never drives them.
This is a simplification with respect to the real world, given the small amount of computational resources available.



\subsection{Salinity Field}